\begin{helper}
Let \(\eta\in\mathbb R\) and \(n\in\mathbb N\).  Put
\[
L=\log n,\quad
k_R=(1+\eta)\sqrt{nL},\quad
K=\noderef{TwoBiteNaturalIndependenceScale}(1+\eta,n),\quad k=(K:\mathbb R),
\]
and
\[
t_1=\frac{\sqrt{nL}}{\log L}.
\]
Assume \(0<\eta\), \(0<n\), \(0<L\), \(0<\log L\), \(0\le k\),
\(k_R\le k\), \(0<t_1\),
\[
\frac{2k}{t_1}+1\le5(1+\eta)\log L,
\]
and
\[
\frac{200(5(1+\eta))^2(\log L)^2L^3}
{(1+\eta)\sqrt{nL}}\le1.
\]
Then
\[
\noderef{RealChooseTwo}\left(\frac{2k}{t_1}+1\right)100L^3
\le\frac12k.
\]
\end{helper}
\begin{proof}
Set
\[
x=\frac{2k}{t_1}+1,\qquad A=5(1+\eta)\log L.
\]
The assumptions \(0\le k\) and \(0<t_1\) imply \(x\ge0\).  By the upper
bound in \(\noderef{RealChooseTwoQuadraticBounds}\),
\[
\noderef{RealChooseTwo}(x)\le x^2.
\]
The displayed \(D\)-estimate gives \(x\le A\), and \(A\ge0\), so
\[
\noderef{RealChooseTwo}(x)\le x^2\le A^2.
\]
Multiplying by the nonnegative quantity \(100L^3\) gives
\[
\noderef{RealChooseTwo}(x)100L^3\le 100A^2L^3.
\]
The threshold assumption has positive denominator and is therefore equivalent
to
\[
200(5(1+\eta))^2(\log L)^2L^3\le k_R.
\]
Since \(A^2=(5(1+\eta))^2(\log L)^2\), the previous inequality yields
\[
100A^2L^3\le\frac12 k_R.
\]
Finally \(k_R\le k\), so \(\frac12k_R\le\frac12k\).  Combining these
inequalities proves the claim.
\end{proof}
